\paragraph{Conclusion.} We developed a lightweight, interpretable context selection layer for EHR question answering that aims to improve the performance of a frozen, resource-constrained (small) language model by selecting a minimal set of evidence chunks under a fixed budget. Empirically, the learned selector achieves performance comparable to strong semantic RAG baselines, outperforming them in some settings but not consistently dominating across all scenarios. The most important qualitative takeaway is \emph{what the model learns}: instead of relying primarily on semantic similarity, it leverages document \emph{structure} (section identity and section--question interactions) to prioritize clinically meaningful regions of the note while down-weighting common distractor sections. This produces a more interpretable and  coherent representation of the EHRs, suggesting a structural redundancy in medical RAG pipelines that can be exploited.

\paragraph{Future work.}
The main bottleneck is supervision quality. Our current chunk labels are weak proxies derived from answer overlap, and teacher-generated QA may introduce bias or omissions; larger gains likely require stronger clinician-verified QA and, ideally, evidence-linked annotations that identify which spans truly support an answer. On the modeling side, richer temporal features (e.g., trends, change-points, event-type-specific windows) and stronger rankers may better capture longitudinal structure and yield more efficient yet relevant contexts. Overall, this project suggests that learning explicit structural representations of clinical notes is a promising direction for improving evidence selection without relying on frontier-scale models at inference time.